% ================= SEKCJA 4 =================
\section{Analiza wymagań}

Analiza wymagań stanowi jeden z kluczowych etapów projektowania systemu informatycznego. Jej celem jest określenie, jakie funkcje powinien realizować system \texttt{WorkFlow}, jakie ograniczenia techniczne należy uwzględnić oraz w jaki sposób poszczególne moduły aplikacji powinny ze sobą współpracować.

W przypadku projektu \texttt{WorkFlow} wymagania wynikają z potrzeby uporządkowania procesów związanych z obsługą pracowników, czasu pracy, dokumentacji kadrowej oraz danych wykorzystywanych przy rozliczeniach. System nie jest projektowany jako pełny system ERP, lecz jako aplikacja wspierająca wybrany, jasno określony zakres procesów kadrowo-organizacyjnych.

\subsection{Podejście do analizy}

Wymagania zostały opracowane z podziałem na trzy główne grupy: wymagania niefunkcyjne, funkcjonalne oraz integracyjne. Taki podział pozwala oddzielić opis konkretnych funkcji systemu od cech jakościowych i technicznych, które wpływają na sposób działania aplikacji.

W analizie przyjęto, że system powinien być możliwy do uruchomienia w środowisku kontenerowym oraz składać się z kilku niezależnych usług. Pozwala to na czytelny podział odpowiedzialności pomiędzy warstwę interfejsu użytkownika, logikę biznesową, bazę danych, magazyn plików oraz reverse proxy.

Z punktu widzenia użytkownika końcowego najważniejsze jest to, aby system umożliwiał sprawną obsługę podstawowych operacji: zarządzanie pracownikami, przypisywanie ich do zakładów i projektów, rejestrowanie czasu pracy, kontrolowanie certyfikatów, przechowywanie dokumentów, obsługę nieobecności oraz przygotowanie danych rozliczeniowych.

\subsection{Zakres wymagań}

Zakres wymagań został dostosowany do aktualnego stanu projektu oraz funkcji przewidzianych do prezentacji. Obejmuje on następujące obszary:
\begin{itemize}
    \item obsługę kont użytkowników i ról systemowych,
    \item zarządzanie pracownikami oraz ich przypisaniem do zakładów,
    \item zarządzanie projektami i przypisaniami pracowników,
    \item rejestrowanie zdarzeń czasu pracy z czytników RCP,
    \item tworzenie i weryfikację wpisów czasu pracy,
    \item obsługę umów pracowników,
    \item obsługę certyfikatów, badań i szkoleń,
    \item rejestrowanie nieobecności,
    \item przechowywanie dokumentów pracowniczych,
    \item przygotowanie danych do rozliczeń,
    \item uruchomienie systemu w środowisku Docker Compose.
\end{itemize}

Z zakresu projektu wyłączono moduły, które nie są częścią aktualnej wersji aplikacji. Dotyczy to w szczególności funkcji związanych z obsługą wypożyczonego sprzętu. Pozwala to zachować spójność pomiędzy dokumentacją, modelem danych prezentowanym w diagramach oraz rzeczywistym zakresem systemu.

\subsection{Założenia techniczne}

Podczas definiowania wymagań przyjęto kilka założeń technicznych, które mają wpływ na dalszy projekt systemu:
\begin{itemize}
    \item aplikacja działa jako system webowy dostępny z poziomu przeglądarki internetowej,
    \item frontend i backend są oddzielnymi komponentami,
    \item backend udostępnia API wykorzystywane przez frontend oraz przez techniczny endpoint czytników RCP,
    \item dane relacyjne są przechowywane w bazie PostgreSQL,
    \item pliki dokumentów są przechowywane poza bazą danych, w magazynie obiektowym MinIO,
    \item komunikacja użytkownika z systemem odbywa się przez Nginx pełniący rolę reverse proxy,
    \item konfiguracja środowiska jest oparta o plik \texttt{.env},
    \item projekt posiada dane demonstracyjne umożliwiające szybkie zaprezentowanie działania systemu.
\end{itemize}

Przyjęta architektura ułatwia rozdzielenie odpowiedzialności pomiędzy poszczególne części systemu. Warstwa frontendowa odpowiada za prezentację danych, warstwa backendowa za logikę biznesową i autoryzację, baza danych za trwałe przechowywanie informacji, a magazyn plików za obsługę dokumentów.

\subsection{Podział wymagań}

Dalsza część rozdziału została podzielona na trzy podrozdziały:
\begin{enumerate}
    \item \textbf{Wymagania niefunkcyjne} -- opisujące cechy jakościowe systemu, takie jak bezpieczeństwo, kontrola dostępu, przenośność środowiska oraz spójność danych.
    \item \textbf{Wymagania funkcjonalne} -- opisujące konkretne operacje wykonywane przez użytkowników oraz główne moduły aplikacji.
    \item \textbf{Wymagania integracyjne} -- opisujące sposób współpracy systemu z bazą danych, magazynem plików, czytnikami RCP oraz usługami infrastrukturalnymi.
\end{enumerate}

Taki układ pozwala uporządkować opis systemu i ułatwia powiązanie wymagań z późniejszym modelem danych, przypadkami użycia oraz scenariuszami działania.

% ================= ROZDZIAŁ 4.1 =================
\subsection{Wymagania niefunkcyjne}

Wymagania niefunkcyjne określają cechy jakościowe systemu \texttt{WorkFlow}, które wpływają na sposób jego działania, bezpieczeństwo, utrzymanie oraz możliwość uruchomienia w przygotowanym środowisku. Nie opisują one pojedynczych funkcji użytkownika, lecz warunki, jakie powinny zostać spełnione przez całą aplikację.

W przypadku systemu \texttt{WorkFlow} szczególne znaczenie mają: bezpieczeństwo danych pracowników, kontrola dostępu, spójność danych, możliwość uruchomienia systemu w środowisku kontenerowym oraz czytelność interfejsu użytkownika.

\subsubsection{Bezpieczeństwo danych}

System przetwarza dane pracowników, informacje o umowach, certyfikatach, nieobecnościach, dokumentach oraz czasie pracy. Z tego względu dostęp do aplikacji powinien być ograniczony do zalogowanych użytkowników.

Podstawowe wymagania bezpieczeństwa obejmują:
\begin{itemize}
    \item uwierzytelnianie użytkowników za pomocą adresu e-mail i hasła,
    \item przechowywanie haseł w postaci zahashowanej,
    \item stosowanie tokenów JWT do autoryzacji żądań,
    \item ograniczenie dostępu do funkcji systemu na podstawie roli użytkownika,
    \item ograniczenie dostępu do danych zakładów zgodnie z przypisanymi uprawnieniami,
    \item zabezpieczenie technicznego endpointu czytników RCP osobnym tokenem urządzenia.
\end{itemize}

W środowisku produkcyjnym system powinien być uruchamiany za pośrednictwem połączenia szyfrowanego HTTPS. W środowisku lokalnym dopuszczalne jest uruchomienie po protokole HTTP, co zostało uwzględnione w konfiguracji plików środowiskowych.

\subsubsection{Kontrola dostępu i role użytkowników}

System powinien realizować kontrolę dostępu opartą na rolach użytkowników. Każda rola odpowiada innemu zakresowi odpowiedzialności w systemie.

W projekcie przewidziano następujące role:
\begin{itemize}
    \item \textbf{Administrator} -- posiada najszerszy zakres dostępu do systemu,
    \item \textbf{HR} -- odpowiada za obsługę danych kadrowych, pracowników, umów, certyfikatów, dokumentów i nieobecności,
    \item \textbf{Księgowość} -- korzysta z danych potrzebnych do rozliczeń,
    \item \textbf{Brygadzista} -- ma dostęp do danych związanych z pracownikami, projektami i zakładem, za który odpowiada,
    \item \textbf{Pracownik} -- reprezentuje podstawowego użytkownika systemu,
    \item \textbf{Użytkownik biurowy} -- rola pomocnicza przewidziana dla prac administracyjnych.
\end{itemize}

Takie rozdzielenie dostępu pozwala ograniczyć widoczność danych wrażliwych oraz zmniejsza ryzyko przypadkowej modyfikacji informacji przez użytkowników, którzy nie powinni mieć dostępu do danego obszaru systemu.

\subsubsection{Spójność i integralność danych}

Dane systemu powinny być przechowywane w relacyjnej bazie danych w sposób zapewniający spójność pomiędzy encjami. Pracownik jest powiązany z zakładem, projektami, umowami, certyfikatami, dokumentami, nieobecnościami oraz wpisami czasu pracy. Z tego względu istotne jest zachowanie poprawnych relacji pomiędzy tabelami.

Wymagania w tym obszarze obejmują:
\begin{itemize}
    \item stosowanie jednoznacznych identyfikatorów rekordów,
    \item wykorzystywanie kluczy obcych do opisu relacji pomiędzy tabelami,
    \item ograniczenie duplikowania danych,
    \item przechowywanie historii wybranych informacji, na przykład umów i przypisań do projektów,
    \item oddzielenie surowych zdarzeń czasu pracy od przetworzonych wpisów czasu pracy.
\end{itemize}

Szczególnie ważne jest rozróżnienie pomiędzy tabelą \texttt{TimeEvent} a tabelą \texttt{TimeEntry}. Pierwsza z nich przechowuje surowe zdarzenia z czytników RCP, natomiast druga zawiera przetworzone wpisy czasu pracy utworzone na podstawie pary zdarzeń wejścia i wyjścia.

\subsubsection{Przenośność środowiska}

System powinien być możliwy do uruchomienia w powtarzalny sposób na różnych środowiskach. W tym celu wykorzystano konteneryzację oraz plik Docker Compose, który definiuje zestaw usług potrzebnych do działania aplikacji.

Środowisko uruchomieniowe powinno obejmować:
\begin{itemize}
    \item kontener frontendu,
    \item kontener backendu,
    \item kontener bazy danych PostgreSQL,
    \item kontener magazynu plików MinIO,
    \item kontener Nginx pełniący funkcję reverse proxy.
\end{itemize}

Takie podejście upraszcza uruchomienie projektu oraz ogranicza liczbę ręcznych kroków konfiguracyjnych. Użytkownik nie musi oddzielnie instalować i konfigurować każdej usługi, ponieważ ich podstawowa konfiguracja znajduje się w pliku \texttt{docker-compose.prod.yml} oraz pliku środowiskowym \texttt{.env}.

\subsubsection{Konfigurowalność}

System powinien umożliwiać zmianę podstawowych parametrów środowiska bez modyfikowania kodu źródłowego. W tym celu wykorzystano zmienne środowiskowe zapisane w pliku \texttt{.env}.

Do najważniejszych parametrów konfiguracyjnych należą:
\begin{itemize}
    \item dane dostępowe do bazy PostgreSQL,
    \item sekret JWT,
    \item token techniczny urządzeń RCP,
    \item dane dostępowe do MinIO,
    \item adres API wykorzystywany przez frontend,
    \item porty usług,
    \item ustawienia związane z ciasteczkami autoryzacyjnymi,
    \item parametry seedów inicjalizujących dane.
\end{itemize}

Do repozytorium nie powinien trafiać rzeczywisty plik \texttt{.env}. Zamiast niego przechowywany jest plik \texttt{.env.example}, który opisuje wymagane zmienne bez ujawniania haseł i tokenów.

\subsubsection{Użyteczność interfejsu}

Interfejs użytkownika powinien być czytelny i dostosowany do typowych zadań administracyjnych. Użytkownik powinien mieć możliwość szybkiego przejścia do najważniejszych modułów, takich jak pracownicy, zakłady, projekty, czas pracy, certyfikaty czy nieobecności.

Wymagania dotyczące użyteczności obejmują:
\begin{itemize}
    \item przejrzysty układ nawigacji,
    \item możliwość wyszukiwania i filtrowania danych,
    \item czytelne formularze dodawania i edycji rekordów,
    \item spójny wygląd widoków w trybie jasnym i ciemnym,
    \item wyraźne oznaczanie statusów, na przykład aktywności pracownika, ważności certyfikatu lub statusu wpisu czasu pracy.
\end{itemize}

System powinien działać w aktualnych wersjach popularnych przeglądarek internetowych. Ponieważ aplikacja ma charakter administracyjny, podstawowym środowiskiem pracy jest komputer z przeglądarką internetową.

\subsubsection{Niezawodność i odtwarzalność danych demonstracyjnych}

Aplikacja powinna umożliwiać ponowne uruchomienie usług bez konieczności ręcznego odtwarzania całej konfiguracji. Dane bazy i magazynu plików powinny być przechowywane w wolumenach Docker, dzięki czemu nie są usuwane przy zwykłym restarcie kontenerów.

Na potrzeby prezentacji projektu przygotowano również dane demonstracyjne. Seed demonstracyjny pozwala odtworzyć przykładowy stan systemu zawierający użytkowników, zakłady, pracowników, projekty, umowy, certyfikaty, nieobecności, dokumenty oraz wpisy czasu pracy. Dzięki temu możliwe jest szybkie przygotowanie środowiska pokazowego bez ręcznego wprowadzania wszystkich danych.
% ================= ROZDZIAŁ 4.2 =================
\subsection{Wymagania funkcjonalne}

Wymagania funkcjonalne określają operacje, które system \texttt{WorkFlow} powinien udostępniać użytkownikom oraz procesy, które powinny być obsługiwane przez aplikację. Wymagania zostały pogrupowane według głównych modułów systemu, aby zachować zgodność pomiędzy opisem funkcjonalnym, przypadkami użycia, modelem danych oraz implementacją.

Zakres funkcjonalny systemu obejmuje obsługę pracowników, zakładów, projektów, czasu pracy, umów, certyfikatów, nieobecności, dokumentów oraz danych rozliczeniowych. Dodatkowo system powinien zapewniać logowanie użytkowników, kontrolę dostępu oraz możliwość uruchomienia danych demonstracyjnych potrzebnych do prezentacji projektu.

\subsubsection{Logowanie i obsługa użytkowników}

System powinien umożliwiać użytkownikowi zalogowanie się do aplikacji przy użyciu adresu e-mail oraz hasła. Po poprawnym logowaniu użytkownik otrzymuje dostęp do panelu systemu zgodnie ze swoją rolą oraz zakresem dostępu do zakładów.

Wymagania w tym obszarze obejmują:
\begin{itemize}
    \item logowanie użytkownika za pomocą adresu e-mail i hasła,
    \item wylogowanie użytkownika z systemu,
    \item przechowywanie informacji o zalogowanym użytkowniku w ramach sesji,
    \item rozpoznawanie roli użytkownika,
    \item ograniczanie dostępu do wybranych widoków i operacji na podstawie roli,
    \item obsługę dostępu użytkownika do jednego lub wielu zakładów.
\end{itemize}

Mechanizm logowania stanowi punkt wejścia do systemu i jest powiązany z modelem pracownika, ponieważ użytkownik aplikacji jest jednocześnie rekordem w tabeli pracowników.

\subsubsection{Zarządzanie zakładami}

System powinien umożliwiać prowadzenie ewidencji zakładów, które reprezentują jednostki organizacyjne firmy. Zakład jest nadrzędnym elementem struktury organizacyjnej, do którego mogą być przypisani pracownicy, projekty oraz czytniki RCP.

Wymagania funkcjonalne dla modułu zakładów obejmują:
\begin{itemize}
    \item dodawanie nowych zakładów,
    \item edycję danych zakładu,
    \item wyświetlanie listy zakładów,
    \item oznaczanie zakładu jako aktywny lub nieaktywny,
    \item przypisywanie pracowników do zakładu,
    \item przypisywanie projektów do zakładu,
    \item obsługę dostępu użytkowników do wybranych zakładów.
\end{itemize}

Dzięki temu system może obsługiwać organizację posiadającą więcej niż jedną lokalizację lub jednostkę administracyjną.

\subsubsection{Zarządzanie pracownikami}

Jednym z głównych modułów systemu jest moduł pracowników. Powinien on umożliwiać prowadzenie centralnej ewidencji osób zatrudnionych lub obsługiwanych w systemie.

System powinien umożliwiać:
\begin{itemize}
    \item dodawanie pracowników,
    \item edycję danych pracownika,
    \item przeglądanie listy pracowników,
    \item wyszukiwanie i filtrowanie pracowników,
    \item przejście do szczegółowego profilu pracownika,
    \item oznaczenie pracownika jako aktywnego lub nieaktywnego,
    \item przypisanie pracownika do zakładu,
    \item przypisanie pracownikowi roli systemowej,
    \item zapis identyfikatora karty RFID wykorzystywanej przy rejestracji czasu pracy.
\end{itemize}

Profil pracownika powinien gromadzić informacje powiązane z daną osobą. Obejmuje to dane podstawowe, umowy, certyfikaty, dokumenty, nieobecności, przypisania do projektów oraz wpisy czasu pracy. Taki układ pozwala traktować profil pracownika jako centralne miejsce dostępu do danych kadrowych.

\subsubsection{Zarządzanie projektami}

System powinien umożliwiać prowadzenie ewidencji projektów. Projekt reprezentuje miejsce, kontrakt lub obszar pracy, do którego mogą być przypisywani pracownicy i z którym mogą być powiązane wpisy czasu pracy.

Wymagania w zakresie projektów obejmują:
\begin{itemize}
    \item tworzenie projektów,
    \item przypisywanie projektu do zakładu,
    \item określanie nazwy, kodu wewnętrznego oraz adresu projektu,
    \item oznaczanie statusu projektu,
    \item określanie daty rozpoczęcia i zakończenia,
    \item przypisywanie pracowników do projektu,
    \item przechowywanie historii przypisań pracowników do projektów.
\end{itemize}

Status projektu powinien umożliwiać rozróżnienie projektów planowanych, aktywnych, zakończonych oraz wstrzymanych. Informacja ta jest istotna przy filtrowaniu projektów oraz przy analizie czasu pracy.

\subsubsection{Przypisania pracowników do projektów}

System powinien umożliwiać powiązanie pracownika z projektem. Przypisanie powinno zawierać datę rozpoczęcia, opcjonalną datę zakończenia oraz dodatkowe uwagi.

Wymagania w tym zakresie obejmują:
\begin{itemize}
    \item przypisanie pracownika do projektu,
    \item zapis daty przypisania,
    \item możliwość zakończenia przypisania,
    \item przechowywanie historii przypisań,
    \item prezentację przypisań w profilu pracownika lub widoku projektu.
\end{itemize}

Mechanizm przypisań pozwala odwzorować fakt, że jeden pracownik może pracować przy różnych projektach w różnych okresach.

\subsubsection{Rejestracja zdarzeń czasu pracy}

System powinien przyjmować zdarzenia czasu pracy pochodzące z czytników RCP. Czytnik jest powiązany z zakładem oraz projektem, dzięki czemu zdarzenie odbicia karty może zostać przypisane do właściwego miejsca pracy.

Zdarzenie czasu pracy powinno zawierać:
\begin{itemize}
    \item numer seryjny czytnika,
    \item identyfikator karty RFID pracownika,
    \item typ zdarzenia: \texttt{IN} lub \texttt{OUT},
    \item dokładny czas zdarzenia,
    \item opcjonalny zewnętrzny identyfikator zdarzenia.
\end{itemize}

Po odebraniu zdarzenia system powinien sprawdzić, czy istnieje aktywny czytnik oraz pracownik o podanym identyfikatorze karty RFID. Surowe zdarzenie powinno zostać zapisane w bazie danych jako rekord \texttt{TimeEvent}.

\subsubsection{Tworzenie wpisów czasu pracy}

Na podstawie zdarzeń \texttt{IN} oraz \texttt{OUT} system powinien tworzyć przetworzone wpisy czasu pracy. Wpis czasu pracy powinien zawierać pracownika, projekt, godzinę rozpoczęcia, godzinę zakończenia, wyliczoną liczbę godzin oraz status.

Wymagania w tym obszarze obejmują:
\begin{itemize}
    \item wyszukanie ostatniego zdarzenia \texttt{IN} dla pracownika,
    \item utworzenie wpisu czasu pracy po odebraniu zdarzenia \texttt{OUT},
    \item powiązanie wpisu z projektem wynikającym z czytnika RCP,
    \item wyliczenie liczby przepracowanych godzin,
    \item oznaczenie wpisu statusem początkowym,
    \item możliwość późniejszej weryfikacji lub zatwierdzenia wpisu.
\end{itemize}

Rozdzielenie tabel \texttt{TimeEvent} i \texttt{TimeEntry} pozwala zachować zarówno surowe dane z urządzenia, jak i dane przetworzone wykorzystywane później w raportowaniu oraz rozliczeniach.

\subsubsection{Obsługa umów}

System powinien umożliwiać przechowywanie informacji o umowach pracowników. Umowa zawiera typ, kwotę wynagrodzenia lub stawkę, datę rozpoczęcia, opcjonalną datę zakończenia oraz informację, czy jest aktualna.

Wymagania funkcjonalne obejmują:
\begin{itemize}
    \item dodanie umowy pracownika,
    \item edycję danych umowy,
    \item przechowywanie typu umowy,
    \item zapis kwoty wynagrodzenia,
    \item określenie daty rozpoczęcia i zakończenia,
    \item oznaczenie umowy jako aktualnej,
    \item przechowywanie historii wcześniejszych umów.
\end{itemize}

Historia umów jest istotna, ponieważ warunki zatrudnienia mogą zmieniać się w czasie. Zachowanie poprzednich rekordów pozwala analizować dane pracownika w kontekście okresu obowiązywania danej umowy.

\subsubsection{Obsługa certyfikatów, badań i szkoleń}

System powinien umożliwiać obsługę certyfikatów, badań lekarskich, szkoleń BHP, uprawnień UDT oraz innych dokumentów potwierdzających kwalifikacje pracownika.

Moduł powinien obejmować:
\begin{itemize}
    \item zarządzanie słownikiem certyfikatów,
    \item określenie typu certyfikatu,
    \item zapis domyślnego okresu ważności,
    \item przypisywanie certyfikatu do pracownika,
    \item zapis numeru certyfikatu lub zaświadczenia,
    \item określenie daty wydania,
    \item określenie daty wygaśnięcia,
    \item prezentację certyfikatów wygasających w wybranym okresie.
\end{itemize}

Funkcja monitorowania dat ważności pozwala szybciej wykryć certyfikaty, które wymagają odnowienia. Jest to szczególnie ważne w przypadku dokumentów wymaganych do wykonywania określonych obowiązków.

\subsubsection{Obsługa nieobecności}

System powinien umożliwiać rejestrowanie nieobecności pracowników. Nieobecność powinna zawierać typ, datę rozpoczęcia, datę zakończenia, status zatwierdzenia oraz opcjonalny dokument.

Wymagania funkcjonalne dla nieobecności obejmują:
\begin{itemize}
    \item dodanie nieobecności pracownika,
    \item określenie typu nieobecności,
    \item wskazanie zakresu dat,
    \item oznaczenie nieobecności jako zatwierdzonej lub oczekującej,
    \item powiązanie nieobecności z dokumentem,
    \item wyświetlanie nieobecności w profilu pracownika.
\end{itemize}

W systemie przewidziano typy nieobecności takie jak urlop, zwolnienie lekarskie, nieobecność nieusprawiedliwiona oraz nieobecność specjalna.

\subsubsection{Obsługa dokumentów}

System powinien umożliwiać przechowywanie informacji o dokumentach pracowników. Plik dokumentu powinien być przechowywany w magazynie obiektowym, natomiast baza danych powinna przechowywać metadane pliku oraz powiązania z innymi rekordami.

Wymagania w tym obszarze obejmują:
\begin{itemize}
    \item zapis nazwy pliku,
    \item zapis lokalizacji pliku w magazynie obiektowym,
    \item powiązanie dokumentu z pracownikiem,
    \item opcjonalne powiązanie dokumentu z umową,
    \item opcjonalne powiązanie dokumentu z certyfikatem,
    \item opcjonalne powiązanie dokumentu z nieobecnością.
\end{itemize}

Takie rozwiązanie pozwala utrzymać logiczne powiązanie dokumentów z danymi kadrowymi bez przechowywania plików bezpośrednio w tabelach relacyjnych.

\subsubsection{Rozliczenia i eksport danych}

System powinien wspierać przygotowanie danych potrzebnych do rozliczeń. Podstawą rozliczeń są zatwierdzone wpisy czasu pracy oraz informacje o pracownikach i ich umowach.

Wymagania funkcjonalne obejmują:
\begin{itemize}
    \item zebranie wpisów czasu pracy z wybranego okresu,
    \item wykorzystanie wyłącznie danych zatwierdzonych,
    \item utworzenie eksportu rozliczeniowego,
    \item zapis miesiąca i roku rozliczeniowego,
    \item zapis informacji o użytkowniku generującym eksport,
    \item powiązanie wpisów czasu pracy z eksportem,
    \item ograniczenie ryzyka ponownego rozliczenia tych samych godzin.
\end{itemize}

Moduł rozliczeniowy nie zastępuje pełnego systemu księgowego. Jego zadaniem jest przygotowanie uporządkowanych danych, które mogą zostać wykorzystane przez osobę odpowiedzialną za rozliczenia.

\subsubsection{Audit log}

System powinien umożliwiać rejestrowanie wybranych operacji wykonywanych na danych. Rejestr audytowy powinien zawierać informację o użytkowniku wykonującym operację, nazwie akcji, encji, identyfikatorze rekordu oraz opcjonalnie poprzednich i nowych wartościach.

Wymagania w tym zakresie obejmują:
\begin{itemize}
    \item zapis użytkownika wykonującego operację,
    \item zapis rodzaju wykonanej akcji,
    \item zapis nazwy encji, której dotyczy operacja,
    \item zapis identyfikatora modyfikowanego rekordu,
    \item możliwość przechowywania wartości przed zmianą i po zmianie.
\end{itemize}

Mechanizm audytu zwiększa przejrzystość pracy systemu i ułatwia analizę zmian w danych wrażliwych, takich jak umowy lub dane pracowników.

\subsubsection{Dane inicjalizacyjne i demonstracyjne}

System powinien posiadać mechanizmy ułatwiające uruchomienie nowego środowiska. W projekcie przewidziano dwa rodzaje danych inicjalizacyjnych.

Pierwszy z nich to seed minimalny, który tworzy podstawowy zakład oraz konta startowe dla głównych ról użytkowników. Drugi to seed demonstracyjny, który generuje pełniejszy zestaw danych potrzebnych do prezentacji działania systemu.

Dane demonstracyjne powinny obejmować:
\begin{itemize}
    \item konta użytkowników,
    \item zakłady,
    \item pracowników,
    \item projekty,
    \item czytniki RCP,
    \item umowy,
    \item certyfikaty,
    \item dokumenty,
    \item nieobecności,
    \item zdarzenia i wpisy czasu pracy,
    \item przykładowy eksport rozliczeniowy.
\end{itemize}

Dzięki temu możliwe jest szybkie przygotowanie środowiska do testów lub prezentacji bez ręcznego wprowadzania danych.
% ================= ROZDZIAŁ 4.3 =================
\subsection{Wymagania integracyjne}

Wymagania integracyjne opisują sposób współpracy systemu \texttt{WorkFlow} z usługami technicznymi, komponentami infrastruktury oraz zewnętrznymi źródłami zdarzeń. W aktualnym zakresie projektu integracje dotyczą przede wszystkim bazy danych, magazynu plików, czytników RCP, komunikacji pomiędzy frontendem i backendem oraz środowiska uruchomieniowego opartego o kontenery.

Celem wymagań integracyjnych jest określenie, jakie elementy muszą ze sobą współpracować, aby system działał jako kompletna aplikacja webowa. Wymagania te nie opisują pojedynczych funkcji użytkownika, lecz zależności pomiędzy usługami i warstwami systemu.

\subsubsection{Integracja z bazą danych}

System powinien integrować się z relacyjną bazą danych PostgreSQL. Baza danych stanowi główne miejsce przechowywania informacji biznesowych, takich jak pracownicy, zakłady, projekty, umowy, certyfikaty, nieobecności, dokumenty, zdarzenia czasu pracy, wpisy czasu pracy oraz eksporty rozliczeniowe.

Komunikacja backendu z bazą danych odbywa się z wykorzystaniem narzędzia Prisma ORM. Schemat bazy został zdefiniowany w pliku \texttt{schema.prisma}, co umożliwia utrzymanie spójności pomiędzy modelem danych aplikacji a strukturą fizycznej bazy danych.

Wymagania w tym obszarze obejmują:
\begin{itemize}
    \item możliwość połączenia backendu z bazą PostgreSQL za pomocą zmiennej środowiskowej,
    \item odwzorowanie encji biznesowych w modelach Prisma,
    \item obsługę relacji pomiędzy tabelami,
    \item możliwość synchronizacji struktury bazy danych ze schematem Prisma,
    \item możliwość inicjalizacji bazy danymi startowymi.
\end{itemize}

\subsubsection{Integracja z magazynem plików}

System powinien współpracować z magazynem obiektowym MinIO. Pliki dokumentów nie są przechowywane bezpośrednio w tabelach bazy danych. W bazie zapisywane są jedynie metadane dokumentu, takie jak nazwa pliku, właściciel dokumentu oraz ścieżka do pliku w magazynie.

Takie rozwiązanie pozwala oddzielić dane relacyjne od plików binarnych. Jest to szczególnie istotne w przypadku dokumentów kadrowych, skanów umów, zaświadczeń, certyfikatów lub dokumentów powiązanych z nieobecnościami.

Wymagania w zakresie integracji z magazynem plików obejmują:
\begin{itemize}
    \item konfigurację endpointu MinIO za pomocą zmiennych środowiskowych,
    \item obsługę danych dostępowych do magazynu plików,
    \item wskazanie nazwy bucketu przechowującego dokumenty,
    \item zapisywanie w bazie danych ścieżki lub adresu pliku,
    \item możliwość powiązania dokumentu z pracownikiem, umową, certyfikatem lub nieobecnością.
\end{itemize}

\subsubsection{Integracja z czytnikami RCP}

System powinien udostępniać techniczny endpoint przeznaczony do odbierania zdarzeń z czytników RCP. Czytnik jest powiązany z zakładem oraz projektem, dzięki czemu zdarzenia wejścia i wyjścia mogą zostać przypisane do konkretnego pracownika oraz miejsca pracy.

Zdarzenie przekazywane do systemu powinno zawierać:
\begin{itemize}
    \item numer seryjny czytnika,
    \item identyfikator karty RFID,
    \item typ zdarzenia: \texttt{IN} albo \texttt{OUT},
    \item czas zdarzenia,
    \item opcjonalny zewnętrzny identyfikator zdarzenia.
\end{itemize}

Endpoint urządzeń powinien być zabezpieczony tokenem technicznym, niezależnym od tokenów użytkowników logujących się do aplikacji. Pozwala to oddzielić komunikację maszynową od standardowej komunikacji użytkownika z systemem.

Po odebraniu zdarzenia backend powinien:
\begin{enumerate}
    \item zweryfikować token urządzenia,
    \item odnaleźć aktywny czytnik na podstawie numeru seryjnego,
    \item odnaleźć aktywnego pracownika na podstawie identyfikatora karty RFID,
    \item zapisać surowe zdarzenie w tabeli \texttt{TimeEvent},
    \item w przypadku zdarzenia \texttt{OUT} spróbować utworzyć wpis czasu pracy na podstawie wcześniejszego zdarzenia \texttt{IN}.
\end{enumerate}

\subsubsection{Komunikacja pomiędzy frontendem i backendem}

Frontend powinien komunikować się z backendem za pomocą API HTTP. Backend udostępnia punkty końcowe odpowiedzialne za logowanie, pobieranie danych, tworzenie rekordów, aktualizację informacji oraz obsługę operacji biznesowych.

Frontend odpowiada za:
\begin{itemize}
    \item prezentację danych użytkownikowi,
    \item obsługę formularzy,
    \item wyświetlanie list i szczegółów rekordów,
    \item kierowanie użytkownika do odpowiednich widoków,
    \item przekazywanie danych do backendu.
\end{itemize}

Backend odpowiada za:
\begin{itemize}
    \item autoryzację żądań,
    \item walidację danych wejściowych,
    \item realizację logiki biznesowej,
    \item komunikację z bazą danych,
    \item komunikację z magazynem plików,
    \item obsługę technicznych endpointów, takich jak odbieranie zdarzeń RCP.
\end{itemize}

Taki podział umożliwia oddzielenie warstwy prezentacji od logiki biznesowej i ułatwia dalsze utrzymanie aplikacji.

\subsubsection{Integracja przez Nginx}

W środowisku Docker Compose system wykorzystuje Nginx jako reverse proxy. Oznacza to, że użytkownik korzysta z jednego adresu aplikacji, a Nginx przekierowuje ruch do odpowiedniej usługi.

Wymagania dotyczące Nginx obejmują:
\begin{itemize}
    \item przekierowanie ruchu aplikacji do frontendu,
    \item przekierowanie zapytań API do backendu,
    \item uproszczenie adresów wykorzystywanych przez użytkownika,
    \item oddzielenie wewnętrznych nazw usług od adresu widocznego z poziomu przeglądarki.
\end{itemize}

Przykładowy podział ruchu wygląda następująco:
\begin{itemize}
    \item ścieżka \texttt{/} kierowana jest do frontendu,
    \item ścieżka \texttt{/api} kierowana jest do backendu.
\end{itemize}

Dzięki temu użytkownik nie musi znać portów poszczególnych usług, a aplikacja może być dostępna pod jednym adresem.

\subsubsection{Środowisko Docker Compose}

System powinien być możliwy do uruchomienia za pomocą Docker Compose. Plik konfiguracyjny opisuje usługi niezbędne do działania aplikacji oraz zależności pomiędzy nimi.

Środowisko powinno obejmować:
\begin{itemize}
    \item usługę frontendu,
    \item usługę backendu,
    \item usługę PostgreSQL,
    \item usługę MinIO,
    \item usługę Nginx.
\end{itemize}

Konfiguracja powinna korzystać ze zmiennych środowiskowych przechowywanych w pliku \texttt{.env}. W repozytorium powinien znajdować się wyłącznie plik przykładowy \texttt{.env.example}, który opisuje wymagane zmienne bez ujawniania rzeczywistych haseł, tokenów i sekretów.

\subsubsection{Inicjalizacja danych}

System powinien umożliwiać przygotowanie bazy danych do pracy po pierwszym uruchomieniu. W projekcie przewidziano dwa mechanizmy inicjalizacji danych:
\begin{itemize}
    \item \textbf{seed minimalny} -- tworzący podstawowy zakład oraz konta startowe dla głównych ról użytkowników,
    \item \textbf{seed demonstracyjny} -- tworzący komplet danych potrzebnych do prezentacji działania systemu.
\end{itemize}

Seed minimalny jest przeznaczony do czystego uruchomienia aplikacji. Seed demonstracyjny służy natomiast do przygotowania środowiska pokazowego i zawiera między innymi pracowników, projekty, czytniki RCP, umowy, certyfikaty, nieobecności, dokumenty, wpisy czasu pracy oraz przykładowy eksport rozliczeniowy.

Ze względu na to, że seed demonstracyjny może czyścić istniejące dane, jego uruchomienie powinno wymagać dodatkowej zmiennej środowiskowej potwierdzającej intencję wykonania operacji.

\subsubsection{Integracje możliwe do rozbudowy}

Aktualna wersja systemu skupia się na podstawowych integracjach wymaganych do działania aplikacji. Przyjęta architektura pozwala jednak na późniejsze rozszerzenie systemu o kolejne mechanizmy.

Do możliwych kierunków rozwoju należą:
\begin{itemize}
    \item integracja z zewnętrznym systemem księgowym,
    \item wysyłka powiadomień e-mail,
    \item integracja z zewnętrznym dostawcą tożsamości,
    \item rozbudowa harmonogramów zadań cyklicznych,
    \item bardziej zaawansowany eksport danych rozliczeniowych.
\end{itemize}

Wymienione elementy nie są wymagane do działania aktualnej wersji projektu. Zostały wskazane jako potencjalne rozszerzenia, które mogą zostać dodane w przyszłości bez zmiany podstawowych założeń architektonicznych systemu.